\documentclass[../../../include/open-logic-section]{subfiles}

\begin{document}

\olfileid[hi]{sth}{z}{milestone}
\olsection{$\Z^-$: एक महत्त्वपूर्ण पड़ाव}

अगले अनुभाग में हम \stagesinf{} पर लौटेंगे। किंतु अनंतता के अभिगृहीत के साथ
हम एक महत्त्वपूर्ण पड़ाव पर पहुँच गए हैं। अब हमारे पास सिद्धांत $\Zminus$
के लिए आवश्यक सभी अभिगृहीत हैं। विस्तार से:

\begin{defn}
सिद्धांत $\Zminus$ के अभिगृहीत ये हैं: विस्तारिता, सम्मिलन, युग्म, घात
समुच्चय, अनंतता, और पृथक्करण योजना के सभी दृष्टांत।
\end{defn}

यह नाम \emph{ज़ेर्मेलो} समुच्चय सिद्धांत को दर्शाता है (उसमें से कुछ
\emph{घटाकर}, जिस पर हम बाद में आएँगे)। ज़ेर्मेलो इस सम्मान के अधिकारी हैं,
क्योंकि उन्होंने मूलतः यह सिद्धांत अपने
\citeyear{Zermelo1908Untersuchungen} में सूत्रित किया था।\footnote{इतिहास और
तकनीकी पक्षों पर रोचक टिप्पणियों के लिए
\citet[Appendix A]{Potter2004} देखें।}

यह सिद्धांत हमें बहुत विशाल मात्रा में गणित करने देने के लिए पर्याप्त
शक्तिशाली है। विशेषतः आपको \olref[sfr][][]{part} में पीछे देखकर स्वयं को
आश्वस्त \emph{करना चाहिए} कि सहज ढंग से किया हमारा हर काम $\Zminus$ के भीतर
अधिक औपचारिक रूप से किया जा सकता है। (कुछ देर ऐसा करने के बाद आप आगे जाकर
\olref[sth][z][arbintersections]{sec} पढ़ना चाह सकते हैं।) अतः आगे, बिना किसी
और टिप्पणी के, हम स्वयं को (कम-से-कम) $\Zminus$ में काम करता मानेंगे।

\end{document}
